% 来源及取材范围见分题文件；此为整理者转录的正文，不是下载原件。
\paragraph{Planar maps.}
\noindent\textit{以下背景与记号据所列原文章节摘编；不是逐字引文，也不是新增证明。}\par
A planar map is a graph endowed with a cyclic permutation of the edges incident to each vertex, such that there exists a proper drawing in which the clockwise order of the curves touching the image of a vertex respects that cyclic permutation.

\begin{theorem}[3.5, The Circle Packing Theorem]
Given any finite simple planar map $G=(V,E)$, $V=\{v_1,\ldots,v_n\}$, there exist $n$ circles in $\R^2$, $C_1,\ldots,C_n$, with disjoint interiors, such that $C_i$ is tangent to $C_j$ if and only if $\{i,j\}\in E$. Furthermore, for every vertex $v_i$, the clockwise order of the circles tangent to $C_i$ agrees with the cyclic permutation of $v_i$'s neighbors in the map.
\end{theorem}
First note that it suffices to prove the theorem for triangulations, that is, simple planar maps in which every face has precisely three edges. Indeed, in any planar map we may add a single vertex inside each face and connect it to all vertices bounding that face. The obtained map is a triangulation, and after applying the circle packing theorem for triangulations, we may remove the circles corresponding to the added vertices, obtaining a circle packing of the original map which respects its cyclic permutations.

\begin{theorem}[3.6]
Let $G=(V,E)$ be a finite triangulation on vertex set $V=\{v_1,\ldots,v_n\}$ and assume that $\{v_1,v_2,v_3\}$ form a face. Then for any three positive numbers $\rho_1,\rho_2,\rho_3$, there exists a circle packing $C_1,\ldots,C_n$ as in Theorem 3.5 with the additional property that $C_1,C_2,C_3$ are mutually tangent, form the outer face, and have radii $\rho_1,\rho_2,\rho_3$, respectively. Furthermore, this circle packing is unique, up to translations and rotations of the plane.
\end{theorem}
\paragraph{Proof of Theorem 3.6: setup.}
We prove Theorem 3.6 which implies Theorem 3.5 as explained above. Therefore we assume from now on that our map is a triangulation. Denote by $n,m,f$ the number of vertices, edges and faces of the map respectively, and observe that $3f=2m$ since each edge is counted in exactly two faces, and each face is bounded by exactly three edges. Therefore, by Euler's formula (Theorem 3.4), we have that
\[
2=n-m+f=n-\tfrac32 f+f=n-\tfrac12f,\qquad f=2n-4.\tag{3.1}
\]
We assume the vertex set is $\{v_1,\ldots,v_n\}$, that $\{v_1,v_2,v_3\}$ is the outer face and that $\rho_1,\rho_2,\rho_3$ are three positive numbers that will be the radii of the outer circles $C_1,C_2,C_3$ eventually. Denote by $F^\circ$ the set of inner faces of the map, and for a subset of vertices $A$ let $F(A)$ be the set of inner faces with at least one vertex in $A$. We write $F(v)$ when we mean $F(\{v\})$.
Given a vector $\mathbf r=(r_1,\ldots,r_n)\in(0,\infty)^n$, an inner face $f\in F^\circ$ bounded by the vertices $v_i,v_j,v_k$, and a distinguished vertex $v_j$, we associate a number $\alpha_f^{\mathbf r}(v_j)=\measuredangle v_iv_jv_k\in(0,\pi)$ which is the angle of $v_j$ in the triangle created by connecting the centers of three mutually tangent circles of radii $r_i,r_j,r_k$ (that is, in a triangle with side lengths $r_i+r_j,r_j+r_k,r_k+r_i$). This number can be calculated using the cosine formula
\[
\cos(\measuredangle v_iv_jv_k)=1-\frac{2r_ir_k}{(r_i+r_j)(r_j+r_k)},
\]
however, we will not use this formula directly. For every $j\in\{1,\ldots,n\}$ we define
\[
\sigma_{\mathbf r}(v_j)=\sum_{f\in F(v_j)}\alpha_f^{\mathbf r}(v_j)
\]
to be the sum of angles at $v_i$ with respect to $\mathbf r$. Let $\theta_1,\theta_2,\theta_3$ be the angles formed at the centers of three mutually tangent circles of radii $\rho_1,\rho_2,\rho_3$. Equivalently, these are the angles of a triangle with edge lengths $r_1+r_2,r_2+r_3$ and $r_1+r_3$. If the vector $\mathbf r$ was the vector of radii of a circle packing of the map satisfying Theorem 3.6, then it would hold that
\[
\sigma_{\mathbf r}(v_i)=\begin{cases}\theta_i&i\in\{1,2,3\},\\2\pi&\text{otherwise},\end{cases}\tag{3.2}
\]
and additionally $(r_1,r_2,r_3)=(\rho_1,\rho_2,\rho_3)$. The proof is split into three parts: finding a positive vector satisfying (3.2); drawing the packing with these radii; and uniqueness up to scaling.

\paragraph{Step 1: Finding the radii vector.}
\begin{claim}[Observation 3.7]
For every $\mathbf r$, $\sum_{i=1}^n\sigma_{\mathbf r}(v_i)=|F^\circ|\pi=(2n-5)\pi$.
\end{claim}
\begin{proof}
Follows immediately since each inner face $f$ bounded by the vertices $v_i,v_j,v_k$ contributes the three angles $\alpha_f^{\mathbf r}(v_i),\alpha_f^{\mathbf r}(v_j),\alpha_f^{\mathbf r}(v_k)$ which sum to $\pi$. By (3.1), there are $2n-5$ inner faces.
\end{proof}
We now set
\[
\delta_{\mathbf r}(v_i)=\begin{cases}\sigma_{\mathbf r}(v_i)-\theta_i&j\in\{1,2,3\},\\\sigma_{\mathbf r}(v_j)-2\pi&\text{otherwise}.\end{cases}\tag{3.3}
\]
Using this notation, our goal is to find $\mathbf r$ for which $\delta_{\mathbf r}\equiv0$. It follows from Observation 3.7 that for every $\mathbf r$,
\[
\sum_{i=1}^n\delta_{\mathbf r}(v_i)=\sum_{i=1}^n\sigma_{\mathbf r}(v_i)-\theta_1-\theta_2-\theta_3-(n-3)2\pi=0.\tag{3.4}
\]
We define $\mathcal E_{\mathbf r}=\sum_{i=1}^n\delta_{\mathbf r}(v_i)^2$. We would like to find $\mathbf r$ for which $\mathcal E_{\mathbf r}=0$.
\begin{claim}[Observation 3.8]
Let $\mathbf r=(r_1,\ldots,r_n)$ and $\mathbf r'=(r_1',\ldots,r_n')$, and let $f\in F^\circ$ be bounded by $v_i,v_j,v_k$.
\begin{itemize}
\item If $r_i'\le r_i,r_k'\le r_k,r_j'\ge r_j$, then $\alpha_f^{\mathbf r'}(v_j)\le\alpha_f^{\mathbf r}(v_j)$.
\item If $r_i'\ge r_i,r_k'\ge r_k,r_j'\le r_j$, then $\alpha_f^{\mathbf r'}(v_j)\ge\alpha_f^{\mathbf r}(v_j)$.
\item $\alpha_f^{\mathbf r}(v_j)$ is continuous in $\mathbf r$.
\end{itemize}
\end{claim}
A proof using the cosine formula is routine and is omitted.

We now define an iterative algorithm, whose input and output are both vectors of radii normalized to have $\ell_1$ norm $1$. We start with the vector $\mathbf r^{(0)}=(1/n,\ldots,1/n)$, and given $\mathbf r=\mathbf r^{(t)}$ we construct $\mathbf r'=\mathbf r^{(t+1)}$. Write $\delta=\delta_{\mathbf r}$ and $\delta'=\delta_{\mathbf r'}$, and similarly $\mathcal E=\mathcal E_{\mathbf r}$ and $\mathcal E'=\mathcal E_{\mathbf r'}$. We begin by ordering the set of reals $\{\delta(v_i)\mid1\le i\le n\}$. If $\delta\equiv0$ we are done; otherwise, we may choose $s\in\R$ such that the set $S=\{v\mid\delta(v)>s\}\ne\varnothing$ and its complement $V\setminus S$ are non-empty and such that the gap
\[
\operatorname{gap}_{\delta}(S):=\min_{v\in S}\delta(v)-\max_{v\notin S}\delta(v)>0
\]
is maximal over all such $s$. Once we choose $S$, a step of the algorithm consists of two steps. For some $\lambda\in(0,1)$ to be chosen later, we set
\[
(\mathbf r_\lambda)_i=\begin{cases}r_i&v_i\in S,\\\lambda r_i&v_i\notin S.\end{cases}
\]
We normalize $\mathbf r_\lambda$ so that the sum of entries is $1$, letting $\bar{\mathbf r}_\lambda$ be the normalized vector. Note that this step does not change the vector $\delta$. We will choose an appropriate $\lambda$ that will decrease all values of $\delta(v)$ for $v\in S$, increase all values of $\delta(v)$ for $v\notin S$, and will close the gap.

\begin{claim}[3.9]
For every $\lambda\in(0,1)$, setting $\mathbf r'=\bar{\mathbf r}_\lambda$, we have that $\delta'(v)\le\delta(v)$ for any $v\in S$, and $\delta'(v)\ge\delta(v)$ for any $v\notin S$.
\end{claim}
\begin{claim}[3.10]
There exists $\lambda\in(0,1)$ such that setting $\mathbf r'=\bar{\mathbf r}_\lambda$ gives that $\operatorname{gap}_{\delta'}(S)=0$.
\end{claim}
\begin{proof}[Proof of Claim 3.9]
Consider $v_j\notin S$ and an inner face $v_i,v_j,v_k$.
\emph{Case I:} $v_i,v_k\notin S$. In this case, the radii of $C_i,C_j,C_k$ are all multiplied by the same number $\lambda$, so $\alpha_f^{\mathbf r'}(v_j)=\alpha_f^{\mathbf r}(v_j)$.
\emph{Case II:} $v_i,v_k\in S$. In this case, the radii of $C_i,C_k$ remain unchanged and the radius of $C_j$ decreases, thus by Observation 3.8, $\alpha_f^{\mathbf r'}(v_j)\ge\alpha_f^{\mathbf r}(v_j)$.
\emph{Case III:} $v_i\notin S,v_k\in S$. In this case the radii of $C_i,C_j$ are multiplied by $\lambda$ and the radius of $C_k$ is unchanged. The angles of the triangle remain unchanged if we multiply all radii by $\lambda^{-1}$, thus we could just as easily have left $C_i,C_j$ unchanged and increased the radius of $C_k$. By Observation 3.8, we get that $\alpha_f^{\mathbf r'}(v_j)\ge\alpha_f^{\mathbf r}(v_j)$.
It follows that $\delta'(v)\ge\delta(v)$ for any $v\notin S$. An identical argument shows that $\delta'(v)\le\delta(v)$ for all $v\in S$.
\end{proof}
\begin{claim}[3.11]
$\displaystyle\lim_{\lambda\downarrow0}\sum_{v\notin S}\delta_{\mathbf r_\lambda}(v)>0$.
\end{claim}
\begin{proof}[Proof of Claim 3.10 using Claim 3.11]
The function $\lambda\mapsto\operatorname{gap}_{\delta_{\mathbf r_\lambda}}(S)$ is continuous on $(0,1]$ by the third bullet of Observation 3.8, and its value at $\lambda=1$ is $\operatorname{gap}_{\delta}(S)>0$. Claim 3.11 says that if $\mu>0$ is small enough, then
$\sum_{v\notin S}\delta_{\mathbf r_\mu}(v)>0$, from which it follows that $\max_{v\notin S}\delta_{\mathbf r_\mu}(v)>0$. By (3.4), we also have $\sum_{v\in S}\delta_{\mathbf r_\mu}(v)<0$, meaning that $\min_{v\in S}\delta_{\mathbf r_\mu}(v)<0$ and therefore $\operatorname{gap}_{\delta_{\mathbf r_\mu}}(S)<0$. By continuity, there exists $\lambda\in(\mu,1)$ such that $\operatorname{gap}_{\delta_{\mathbf r_\lambda}}(S)=0$.
\end{proof}
\begin{proof}[Proof of Claim 3.11]
We first show that for each face $f\in F(V\setminus S)$ bounded by $v_i,v_j,v_k$, the sum of angles at the vertices belonging to $V\setminus S$ converges to $\pi$ as $\lambda\downarrow0$. We show this by the following case analysis. The statements in cases II and III can be justified by drawing a picture or appealing to the cosine formula.
\emph{Case I:} If $v_i,v_j,v_k\notin S$ then since the face is a triangle,
$\alpha_f^{\mathbf r_\lambda}(v_i)+\alpha_f^{\mathbf r_\lambda}(v_j)+\alpha_f^{\mathbf r_\lambda}(v_k)=\pi$ for all $\lambda\in(0,1)$.
\emph{Case II:} If $v_i,v_j\notin S$ but $v_k\in S$ then $\lim_{\lambda\downarrow0}\alpha_f^{\mathbf r_\lambda}(v_k)=0$, hence $\lim_{\lambda\downarrow0}(\alpha_f^{\mathbf r_\lambda}(v_i)+\alpha_f^{\mathbf r_\lambda}(v_j))=\pi$.
\emph{Case III:} If $v_i\notin S$ but $v_j,v_k\in S$ then $\lim_{\lambda\downarrow0}(\alpha_f^{\mathbf r_\lambda}(v_j)+\alpha_f^{\mathbf r_\lambda}(v_k))=0$, hence $\lim_{\lambda\downarrow0}\alpha_f^{\mathbf r_\lambda}(v_i)=\pi$.
It follows that
\[
\lim_{\lambda\downarrow0}\sum_{v\notin S}\sigma_{\mathbf r_\lambda}(v)=|F(V\setminus S)|\pi.\tag{3.5}
\]
For convenience, set $\theta(v_i)=\theta_i$ for $1\le i\le3$, and $\theta(v_i)=2\pi$ otherwise, so that $\delta_{\mathbf r}(v)=\sigma_{\mathbf r}(v)-\theta(v)$ for all $v\in V$. Then
\[
\lim_{\lambda\downarrow0}\sum_{v\notin S}\delta_{\mathbf r_\lambda}(v)
=|F(V\setminus S)|\pi-\sum_{v\notin S}\theta(v).\tag{3.6}
\]
Let $\bar F=F^\circ\setminus F(V\setminus S)$, so every face in $\bar F$ contains only vertices of $S$. We will show that
\[
|\bar F|\pi<\sum_{v\in S}\theta(v).\tag{3.7}
\]
If (3.7) holds, then we can add the negative quantity $|\bar F|\pi-\sum_{v\in S}\theta(v)$ to the right side of (3.6), obtaining $|F^\circ|\pi-\sum_{v\in V}\theta(v)=(2n-5)\pi-(2n-5)\pi=0$. It follows that (3.6) is strictly positive, proving the claim. Thus it suffices to show (3.7).

In the rest of the proof, we fix an embedding of $G$ in the plane with $(v_1,v_2,v_3)$ as the outer face. Let $G[S]$ be the subgraph of $G$ induced by $S$. Partition $S$ into equivalence classes, $S=S_1\cup\cdots\cup S_k$, where two vertices are equivalent if they are in the same connected component of $G[S]$. Then $G[S]=G[S_1]\cup\cdots\cup G[S_k]$. Let $\bar F_j$ be the set of faces in $\bar F$ that appear as faces of $G[S_j]$, so that we have the disjoint union $\bar F=\bar F_1\cup\cdots\cup\bar F_k$. Since $S$ is nonempty, it is enough to show that for all $1\le j\le k$,
\[
|\bar F_j|\pi<\sum_{v\in S_j}\theta(v).\tag{3.8}
\]
Let $m_j$ and $f_j$ denote the number of edges and faces, respectively, of $G[S_j]$. Observe that $|\bar F_j|\le f_j-1$. If $|\bar F_j|=0$, then (3.8) is trivial. If $|\bar F_j|\ge1$, then $G[S_j]$ has at least one inner face, and since it is a simple graph, every face must have degree at least $3$. (The degree of a face is the number of directed edges that make up its boundary.) Because the sum of the degrees of all the faces equals twice the number of edges, we have $2m_j\ge3f_j$. Euler's formula now gives
\[
|S_j|+f_j-2=m_j\ge\tfrac32 f_j,
\]
and hence $f_j\le2|S_j|-4$. Thus, the left side of (3.8) satisfies $|\bar F_j|\pi\le(2|S_j|-5)\pi$.
If $S_j$ contains all of $v_1,v_2,v_3$, then the right side of (3.8) is
$\theta_1+\theta_2+\theta_3+(|S_j|-3)2\pi=(2|S_j|-5)\pi$. Otherwise, at least one of the $\theta_i$ is replaced by $2\pi$ and so the right side of (3.8) is strictly greater than the left side. In fact, (3.8) holds except when $v_1,v_2,v_3\in S_j$ and $|\bar F_j|=f_j-1=2|S_j|-5$. We now show that this situation cannot occur.
The equality $|\bar F_j|=f_j-1$ means that every inner face of $G[S_j]$ is an element of $\bar F_j$ and therefore a face of $G$. Since $v_1,v_2,v_3\in S_j$, the outer face of $G[S_j]$ is $(v_1,v_2,v_3)$, which is the same as the outer face of $G$. So, every face of $G[S_j]$ is also a face of $G$. But this is impossible: if we choose any $v\in V\setminus S$, then $v$ must lie in some face of $G[S_j]$, which then cannot be a face of $G$. Therefore, it cannot be true that $v_1,v_2,v_3\in S_j$ and also $|\bar F_j|=f_j-1$, so we conclude that (3.8) always holds.
\end{proof}

We now analyse the algorithm. Let $\lambda\in(0,1)$ be the one guaranteed by Claim 3.10, and set $\mathbf r'=\bar{\mathbf r}_\lambda$.
\begin{claim}[3.12]
$\mathcal E'\le\mathcal E(1-1/(2n^3))$.
\end{claim}
\begin{proof}
Define $t=\min_{v\in S}\delta'(v)=\max_{v\notin S}\delta'(v)$. By (3.4) we have that $\sum_{i=1}^n\delta(v_i)=\sum_{i=1}^n\delta'(v_i)=0$, hence
\begin{align*}
\mathcal E-\mathcal E'
&=\sum_i\delta(v_i)^2-\sum_i\delta'(v_i)^2\\
&=\sum_i(\delta(v_i)-\delta'(v_i))^2
+2\sum_i(t-\delta'(v_i))(\delta'(v_i)-\delta(v_i)).
\end{align*}
If $v\in S$, then $t\le\delta'(v)\le\delta(v)$ and if $v\notin S$, then $t\ge\delta'(v)\ge\delta(v)$. Thus, in both cases $(t-\delta'(v))(\delta'(v)-\delta(v))\ge0$. Taking $u\in S$ and $v\notin S$ with $\delta'(u)=\delta'(v)=t$, we have that
\[
\mathcal E-\mathcal E'\ge(\delta(u)-t)^2+(\delta(v)-t)^2
\ge\frac{(\delta(u)-\delta(v))^2}{2}\ge\frac{\operatorname{gap}_\delta(S)^2}{2}.
\]
Since $\operatorname{gap}_\delta(S)$ was chosen to be the maximal gap we may bound,
\[
\operatorname{gap}_\delta(S)\ge\frac1n\left(\max_{v\in V}\delta(v)-\min_{v\in V}\delta(v)\right).
\]
For every $v\in V$, $\max_w\delta(w)-\min_w\delta(w)\ge|\delta(v)|$, and thus
\[
n\left(\max_v\delta(v)-\min_v\delta(v)\right)^2\ge\sum_i\delta(v_i)^2=\mathcal E.
\]
Hence $\mathcal E-\mathcal E'\ge\frac1{2n^2}(\max_v\delta(v)-\min_v\delta(v))^2\ge\frac1{2n^2}\frac{\mathcal E}{n}$, and we conclude that $\mathcal E'\le\mathcal E(1-1/(2n^3))$.
\end{proof}
Write $\mathcal E^{(t)}=\mathcal E_{\mathbf r^{(t)}}$. By iterating the described algorithm, we obtain from Claim 3.12 that
\[
\mathcal E^{(t)}\le\mathcal E^{(0)}\left(1-\frac1{2n^3}\right)^t\longrightarrow0\quad\text{as }t\to\infty.
\]
By our normalization $\|\mathbf r^{(t)}\|_{\ell_1}=1$. Thus, by compactness, there exists a subsequence $\{t_k\}$ and a vector $\mathbf r^\infty$ such that $\mathbf r^{(t_k)}\to\mathbf r^\infty$ as $k\to\infty$. From continuity of $\mathcal E$ we have that $\mathcal E(\mathbf r^\infty)=0$, meaning that (3.2) is satisfied. For $\mathbf r^\infty$ to be feasible as a vector of radii, we also have to argue that it is positive (the fact that no coordinates are $\infty$ follows since $\|\mathbf r^\infty\|_{\ell_1}=1$).
\begin{claim}[3.13]
$\mathbf r_i^\infty>0$ for every $i$.
\end{claim}
\begin{proof}
Let $S=\{v_i\in V:\mathbf r_i^\infty>0\}$. Because of the normalization of $\mathbf r$, we know that $S$ is nonempty. Assume for contradiction that $S\subsetneq V$. We repeat the exact same argument used in the proof of Claim 3.11 showing first by case analysis that
\[
\lim_{t\to\infty}\sum_{v\notin S}\sigma_{\mathbf r^{(t)}}(v)=|F(V\setminus S)|\pi
\]
and then deducing that $\lim_{t\to\infty}\sum_{v\notin S}\delta_{\mathbf r^{(t)}}(v)>0$.
This contradicts that $\lim_{t\to\infty}\mathcal E^{(t)}=0$, so we conclude that $S=V$.
\end{proof}

\paragraph{Step 2: Drawing the circle packing described by $\mathbf r^\infty$.}
Given the vector of radii $\mathbf r^\infty$ satisfying (3.2), we now show that the corresponding circle packing can be drawn uniquely up to translations and rotations. In fact, we provide a slightly more general statement which is due to Ori Gurel-Gurevich and Ohad Feldheim [personal communications, 2018].
Let $G=(V,E)$ be a finite planar triangulation on vertex set $\{v_1,\ldots,v_n\}$ and assume that $\{v_1,v_2,v_3\}$ is the outer face. A vector of positive real numbers $\boldsymbol\ell=\{\ell_e\}_{e\in E}$ indexed by the edge set $E$ is called feasible if for any face enclosed by edges $e_1,e_2,e_3$, the lengths $\ell_{e_1},\ell_{e_2},\ell_{e_3}$ can be made to form a triangle. In other words, these lengths satisfy three triangle inequalities,
$\ell_{e_i}+\ell_{e_j}>\ell_{e_k}$ for $\{i,j,k\}=\{1,2,3\}$.
Given a feasible edge length vector $\boldsymbol\ell$ we may again use the cosine formula to compute, for each face $f$, the angle at a vertex of the triangle formed by the three corresponding edge lengths. We denote these angles, as before, by $\alpha_f^{\boldsymbol\ell}(v)$ where $v$ is a vertex of $f$. Similarly, we define $\sigma_{\boldsymbol\ell}(v)=\sum_{f\in F(v)}\alpha_f^{\boldsymbol\ell}(v)$ to be the sum of angles at a vertex $v$.
\begin{theorem}[3.14]
Let $G$ be a finite triangulation and $\boldsymbol\ell$ a feasible vector of edge lengths. Assume that $\sigma_{\boldsymbol\ell}(v)=2\pi$ for any internal vertex $v$. Then there is a drawing of $G$ in the plane so that each edge $e$ is drawn as a straight line segment of length $\ell_e$ and no two edges cross. Furthermore, this drawing is unique up to translations and rotations.
\end{theorem}
It is easy to use the theorem above to draw the circle packing given the radii vector $\mathbf r^\infty$ satisfying (3.2). Indeed, given $\mathbf r^\infty$ we set $\boldsymbol\ell$ by putting $\ell_e=\mathbf r_i^\infty+\mathbf r_j^\infty$ for any edge $e=\{v_i,v_j\}$ of the graph. Condition (3.2) implies that $\boldsymbol\ell$ is feasible. We now apply Theorem 3.14 and obtain the guaranteed drawing and draw a circle $C_i$ of radii $\mathbf r_i^\infty$ around $v_i$ for all $i$.
Theorem 3.14 guarantees that for any edge $\{v_i,v_k\}$ the distance between $v_i,v_j$ is precisely $\mathbf r_i^\infty+\mathbf r_j^\infty$ and thus $C_i$ and $C_j$ are tangent. Conversely, assume that $v_i,v_j$ do not form an edge. To each vertex $v$ let $A_v$ be the union of triangles touching $v$, each triangle is the space bounded by a face touching $v$ in the drawing of Theorem 3.14. Since $G$ is a triangulation and $v_i,v_j$ are not adjacent we learn that $A_{v_i}$ and $A_{v_j}$ have disjoint interiors. Furthermore, $C_i\subset\operatorname{Int}(A_{v_i})$ since the straight lines emanating from $v_i$ have length larger than $\mathbf r_i^\infty$. By the same token $C_j\subset\operatorname{Int}(A_{v_j})$ and we conclude that $C_i$ and $C_j$ are not tangent.
Lastly, we note that by (3.2) the outer boundary of the polygon we drew is a triangle with angles $\theta_1,\theta_2,\theta_3$ and hence $(r_1,r_2,r_3)$ is a positive multiple of $(\rho_1,\rho_2,\rho_3)$. Step 2 of the proof of Theorem 3.5 is now concluded.
\begin{proof}[Proof of Theorem 3.14]
We prove this by induction on the number of vertices $n$. The base case $n=3$ is trivial since the feasibility of $\boldsymbol\ell$ guarantees that the edge lengths of the three edges of the outer face can form a triangle. Any two triangles with the same edge lengths can be rotated and translated to be identical, so the uniqueness statement holds for $n=3$.
Assume now that $n>3$ so that there exists an internal vertex $v$. Denote by $v_1,\ldots,v_m$ the neighbors of $v$ ordered clockwise. We begin by placing $v$ at the origin and drawing all the faces to which $v$ belongs. That is, we draw the edge $\{v,v_1\}$ as a straight line interval of length $\ell_{\{v,v_1\}}$ on the positive $x$-axis emanating from the origin and proceed iteratively: for each $1<i\le m$ we draw the edge $\{v,v_i\}$ as a straight line interval of length $\ell_{\{v,v_i\}}$ emanating from the origin ($v$) at a clockwise angle of $\alpha_f^{\boldsymbol\ell}(v)$ from the previous drawn line segment of $\{v,v_{i-1}\}$, where $f=\{v,v_{i-1},v_i\}$.
This determines the location of $v_1,\ldots,v_m$ in the plane and allows us to ``complete'' the triangles by drawing the straight line segments connecting $v_i$ to $v_{i+1}$, each of length $\ell_{\{v_i,v_{i+1}\}}$ where $1\le i\le m$ (where $v_{m+1}=v_1$). Denote these edges by $e_1,\ldots,e_m$.
Since $\sigma_{\boldsymbol\ell}(v)=2\pi$ we learn that these $m$ triangles have disjoint interiors and that the edges $e_1,\ldots,e_m$ form a closed polygon containing the origin in its interior. It is a classical fact [G.H. Meisters, \emph{Polygons have ears}, 1975] that every closed polygon can be triangulated by drawing some diagonals as straight line segments in the interior of the polygon.
We fix such a choice of diagonals and use it to form a new graph $G'$ on $n-1$ vertices and $|E(G)|-3$ edges by erasing $v$ and the $m$ edges emanating from it and adding the new $m-3$ edges corresponding to the diagonals we added. Furthermore, we generate a new edge length vector $\boldsymbol\ell'$ corresponding to $G'$ by assigning the new edges lengths corresponding to the Euclidean length of the drawn diagonals and leaving the other edge lengths unchanged.
It is clear that $\boldsymbol\ell'$ is feasible and that the angle sum at each internal vertex of $G'$ is $2\pi$. Therefore we may apply the induction hypothesis and draw the graph $G'$ according to the edge lengths $\boldsymbol\ell'$. This drawing is unique up to translations and rotations by induction. Note that in this drawing of $G'$, the polygon corresponding to $e_1,\ldots,e_m$ must be the exact same polygon as before, up to translations and rotations, since it has the same edge lengths and the same angles between its edges.
Since it is the same polygon, we can now erase the diagonals in this drawing and place a new vertex in the same relative location where we drew $v$ previously, along with the straight line segments connecting it to $v_1,\ldots,v_m$. Thus we have obtained the desired drawing of $G$. The uniqueness up to translations and rotations of this drawing follows from the uniqueness of the drawing of $G'$ and the fact that the location of $v$ is uniquely determined in that drawing.
\end{proof}

\paragraph{Step 3: Uniqueness.}
\begin{theorem}[3.15, Uniqueness of Circle Packing]
Given a simple finite triangulation with outer face $v_1,v_2,v_3$ and three radii $\rho_1,\rho_2,\rho_3$, the circle packing with $C_{v_1},C_{v_2},C_{v_3}$ having radii $\rho_1,\rho_2,\rho_3$ is unique up to translations and rotations.
\end{theorem}
\begin{proof}
We have already seen in step 2 that given the radii vector $\mathbf r$ the drawing we obtain is unique up to translations and rotations. Thus, we only need to show the uniqueness of $\mathbf r$ given $\rho_1,\rho_2,\rho_3$.
To that aim, suppose that $\mathbf r^a$ and $\mathbf r^b$ are two vectors satisfying (3.2). Since the outer face in both vectors correspond to a triangle of angles $\theta_1,\theta_2,\theta_3$ we may rescale so that $\mathbf r_i^a=\mathbf r_i^b=\rho_i$ for $i=1,2,3$. After this rescaling, assume by contradiction that $\mathbf r^a\ne\mathbf r^b$ and let $v$ be the interior vertex which maximizes $\mathbf r_v^a/\mathbf r_v^b$. We can assume without loss of generality that this quantity is strictly larger than $1$, as otherwise we can swap $\mathbf r^a$ and $\mathbf r^b$.
Now we claim that for each $f=(v,u_1,u_2)\in F(v)$, we have $\alpha_f^{\mathbf r^a}(v)\le\alpha_f^{\mathbf r^b}(v)$, with equality if and only if the ratios $\mathbf r_{u_i}^a/\mathbf r_{u_i}^b$, for $i=1,2$, are both equal to $\mathbf r_v^a/\mathbf r_v^b$. This is a direct consequence of Observation 3.8.
Indeed, scale all the radii in $\mathbf r^b$ by a factor of $\mathbf r_v^a/\mathbf r_v^b$ to get a new vector $\mathbf r'$ such that $\mathbf r_v^a=\mathbf r_v'$ and $\mathbf r_u^a\le\mathbf r_u'$ for all $u\ne v$. The second bullet point in Observation 3.8 implies that $\alpha_f^{\mathbf r^a}(v)\le\alpha_f^{\mathbf r'}(v)=\alpha_f^{\mathbf r^b}(v)$. As well, if either $\mathbf r_{u_1}^a<\mathbf r_{u_1}'$ or $\mathbf r_{u_2}^a<\mathbf r_{u_2}'$, then the cosine formula yields the strict inequality $\alpha_f^{\mathbf r^a}(v)<\alpha_f^{\mathbf r'}(v)$. Thus, $\alpha_f^{\mathbf r^a}(v)=\alpha_f^{\mathbf r^b}(v)$ only if $\mathbf r_{u_i}^a/\mathbf r_{u_i}^b=\mathbf r_v^a/\mathbf r_v^b$ for $i=1,2$.
Now, since $\alpha_f^{\mathbf r^a}(v)\le\alpha_f^{\mathbf r^b}(v)$ for each $f\in F(v)$, while $\sigma_{\mathbf r^a}(v)=\sigma_{\mathbf r^b}(v)=2\pi$, the equality $\alpha_f^{\mathbf r^a}(v)=\alpha_f^{\mathbf r^b}(v)$ must hold for each $f$. Therefore, each neighbor $u$ of $v$ satisfies $\mathbf r_u^a/\mathbf r_u^b=\mathbf r_v^a/\mathbf r_v^b$.
Because the graph is connected, the ratio $\mathbf r_u^a/\mathbf r_u^b$ must be constant for all vertices $u\in V(G)$. But this contradicts that $\mathbf r_v^a/\mathbf r_v^b>1$ while $\mathbf r_{v_i}^a/\mathbf r_{v_i}^b=1$ for $i=1,2,3$. We conclude that $\mathbf r^a=\mathbf r^b$.
\end{proof}
